\section{Details: Infinite products (part 1)}
\label{sec.details.gf.prod}

Fix a commutative ring $K$.

\subsection{Essentially finite families}

\begin{definition}
\label{def.fps.infprod.essFinite}
\lean{PowerSeries.EssentiallyFinite}
\leanhelper
\leanok
A family $(k_i)_{i \in I}$ of natural numbers is \emph{essentially finite}
if all but finitely many entries equal $0$, i.e., the set $\{i \in I : k_i \neq 0\}$
is finite. This corresponds to $S^I_{\mathrm{fin}}$ in the source.
\end{definition}

\begin{definition}
\label{def.fps.infprod.coeffSupportSet}
\lean{PowerSeries.CoeffSupportSet}
\leanhelper
\leanok
For a family $(p_{i,k})$ of FPSs, the \emph{coefficient support set} $T_m$ at degree $m$
is the set $\{(i,k) : k \neq 0 \text{ and } [x^m] p_{i,k} \neq 0\}$.
\end{definition}

\begin{definition}
\label{def.fps.infprod.coeffSupportSetUnion}
\lean{PowerSeries.CoeffSupportSetUnion}
\leanhelper
\leanok
The \emph{coefficient support union} $T'_n = T_0 \cup T_1 \cup \cdots \cup T_n$
is the union of coefficient support sets up to degree $n$.
\end{definition}

\begin{definition}
\label{def.fps.infprod.indexSetIn}
\lean{PowerSeries.IndexSetIn}
\leanhelper
\leanok
The \emph{index set} $I_n$ is the set of first components appearing in $T'_n$:
$I_n = \{i \in I : \exists k,\; (i,k) \in T'_n\}$.
\end{definition}

\begin{definition}
\label{def.fps.infprod.valueSetKn}
\lean{PowerSeries.ValueSetKn}
\leanhelper
\leanok
The \emph{value set} $K_n$ is the set of second components appearing in $T'_n$:
$K_n = \{k \in \mathbb{N} : \exists i,\; (i,k) \in T'_n\}$.
\end{definition}

\begin{definition}
\label{def.fps.infprod.SfinI}
\lean{PowerSeries.SfinI}
\leanhelper
\leanok
The set $S^I_{\mathrm{fin}}$ of \emph{essentially finite families} in $\prod_{i \in I} S_i$
is the set of all families $(k_i)_{i \in I}$ such that $k_i \in S_i$ for all $i$ and
the family $(k_i)$ is essentially finite.
\end{definition}

\begin{definition}
\label{def.fps.infprod.extendFromFinset}
\lean{PowerSeries.extendFromFinset}
\leanhelper
\leanok
Given a finite subset $I_n \subseteq I$ and a function $f : I_n \to \mathbb{N}$,
\emph{extend from finset} produces a function $I \to \mathbb{N}$ by setting
$f(i)$ for $i \in I_n$ and $0$ otherwise. This is used in the bijection between
essentially finite families supported on $I_n$ and functions $I_n \to S$.
\end{definition}

\begin{definition}
\label{def.fps.infprod.restrictToFinset}
\lean{PowerSeries.restrictToFinset}
\leanhelper
\leanok
Given a finite subset $I_n \subseteq I$ and a function $k : I \to \mathbb{N}$,
\emph{restrict to finset} produces the restriction $k|_{I_n} : I_n \to \mathbb{N}$.
This is the inverse direction of the bijection used in Claim~8.
\end{definition}

\subsection{Existence of approximators}

\begin{lemma}
\label{lem.fps.mulable.approx.lean}
\lean{PowerSeries.exists_xn_approximator}
\leanhelper
\leanok
For a multipliable family $(\mathbf{a}_i)_{i \in I}$ and any $n \in \mathbb{N}$,
there exists an $x^n$-approximator $M$ for $(\mathbf{a}_i)_{i \in I}$.
The proof takes the union of finite sets $M_0, M_1, \ldots, M_n$ where each $M_k$
determines the $x^k$-coefficient.
\end{lemma}

\subsection{Multipliability criteria}

\begin{theorem}
\label{thm.fps.1+f-mulable.lean}
\lean{PowerSeries.multipliable_one_add_of_summable}
\leanhelper
\leanok
If $(f_i)_{i \in I}$ is a family of FPSs such that for each $n$,
the set $\{i \in I : [x^n] f_i \neq 0\}$ is finite (i.e., the family is
coefficient-wise summable), then the family $(1 + f_i)_{i \in I}$ is multipliable.

The proof constructs an approximator $M$ as the union of the finite sets
$\{i : [x^m] f_i \neq 0\}$ for $m = 0, 1, \ldots, n$, and shows that
multiplying by $(1 + f_i)$ for $i \notin M$ does not change coefficients up to degree $n$.
\end{theorem}

\begin{theorem}
\label{thm.fps.1-mulable.lean}
\lean{PowerSeries.multipliable_of_finite_ne_one}
\leanhelper
\leanok
If $(a_i)_{i \in I}$ is a family of FPSs such that
$\{i \in I : a_i \neq 1\}$ is finite, then the family is multipliable.
The finite support set itself serves as a determining set for every coefficient.
\end{theorem}

\begin{theorem}
\label{thm.fps.union-mulable.lean}
\lean{PowerSeries.multipliable_of_union}
\leanhelper
\leanok
If $J \subseteq I$ and both the subfamily $(a_i)_{i \in J}$ and the subfamily
$(a_i)_{i \in I \setminus J}$ are multipliable, then the full family $(a_i)_{i \in I}$
is multipliable. The proof combines $x^n$-approximators for the two subfamilies.
\end{theorem}

\subsection{Multipliable product and division rules (Lean formalization)}

\begin{theorem}
\label{thm.fps.prod-mulable.lean}
\lean{PowerSeries.multipliable_mul}
\leanhelper
\leanok
If $(\mathbf{a}_i)_{i \in I}$ and $(\mathbf{b}_i)_{i \in I}$ are both multipliable,
then the pointwise product family $(\mathbf{a}_i \cdot \mathbf{b}_i)_{i \in I}$ is multipliable.
The proof takes $M = U \cup V$ where $U$ and $V$ are $x^n$-approximators for the two families.
\end{theorem}

\begin{theorem}
\label{thm.fps.prod-mulable-eq.lean}
\lean{PowerSeries.tprod_mul_eq_mul_tprod}
\leanhelper
\leanok
Under the hypotheses of Theorem~\ref{thm.fps.prod-mulable.lean}, the infinite product
of the pointwise product equals the product of the two infinite products:
\[
\prod_{i \in I} (\mathbf{a}_i \cdot \mathbf{b}_i)
= \Bigl(\prod_{i \in I} \mathbf{a}_i\Bigr) \cdot \Bigl(\prod_{i \in I} \mathbf{b}_i\Bigr).
\]
\end{theorem}

\begin{theorem}
\label{thm.fps.div-mulable.lean}
\lean{PowerSeries.multipliable_div}
\leanhelper
\leanok
If $(\mathbf{a}_i)_{i \in I}$ and $(\mathbf{b}_i)_{i \in I}$ are multipliable and
each $\mathbf{b}_i$ has invertible constant term, then the family
$(\mathbf{a}_i \cdot \mathbf{b}_i^{-1})_{i \in I}$ is multipliable.
\end{theorem}

\begin{theorem}
\label{thm.fps.div-mulable-eq.lean}
\lean{PowerSeries.tprod_div_eq_div_tprod}
\leanhelper
\leanok
Under the hypotheses of Theorem~\ref{thm.fps.div-mulable.lean},
\[
\prod_{i \in I} \frac{\mathbf{a}_i}{\mathbf{b}_i}
= \frac{\prod_{i \in I} \mathbf{a}_i}{\prod_{i \in I} \mathbf{b}_i}.
\]
\end{theorem}

\subsection{Subfamily and reindexing lemmas}

\begin{lemma}
\label{lem.fps.infprod.approx-inter-subfamily}
\lean{PowerSeries.isXnApproximator_inter_subfamily}
\leanhelper
\leanok
If $U$ is an $x^n$-approximator for a family $(\mathbf{a}_i)_{i \in I}$ of
invertible FPSs and $J \subseteq I$, then $U \cap J$ (as a preimage under the
subtype embedding) is an $x^n$-approximator for the subfamily
$(\mathbf{a}_i)_{i \in J}$.
This is the Lean formalization of Lemma~\ref{lem.fps.prods-mulable-subfams-appr},
using the invertibility hypothesis to cancel factors outside $J$.
\end{lemma}

\begin{theorem}
\label{thm.fps.prods-mulable-subfams.lean}
\lean{PowerSeries.multipliable_subfamily}
\leanhelper
\leanok
If $(\mathbf{a}_i)_{i \in I}$ is multipliable and each $\mathbf{a}_i$ has invertible
constant term, then for any $J \subseteq I$, the subfamily $(\mathbf{a}_i)_{i \in J}$
is multipliable.
\end{theorem}

\begin{theorem}
\label{thm.fps.prods-mulable-reindex.lean}
\lean{PowerSeries.multipliable_reindex}
\leanhelper
\leanok
If $f : S \to T$ is a bijection and $(\mathbf{a}_t)_{t \in T}$ is multipliable,
then $(\mathbf{a}_{f(s)})_{s \in S}$ is multipliable.
\end{theorem}

\begin{theorem}
\label{thm.fps.prods-mulable-reindex-eq.lean}
\lean{PowerSeries.tprod_reindex}
\leanhelper
\leanok
Under the hypotheses of Theorem~\ref{thm.fps.prods-mulable-reindex.lean},
the reindexed infinite product equals the original:
$\prod_{s \in S} \mathbf{a}_{f(s)} = \prod_{t \in T} \mathbf{a}_t$.
\end{theorem}

\subsection{Fiber product decomposition}

\begin{theorem}
\label{thm.fps.prods-mulable-fibers.lean}
\lean{PowerSeries.multipliable_prod_fibers}
\leanhelper
\leanok
Let $f : S \to W$ be a map, and $(\mathbf{a}_s)_{s \in S}$ a multipliable family.
Suppose that for each $w \in W$, the fiber subfamily
$(\mathbf{a}_s)_{s \in f^{-1}(w)}$ is multipliable.
Then the family of fiber products
$\bigl(\prod_{s \in f^{-1}(w)} \mathbf{a}_s\bigr)_{w \in W}$ is multipliable.
The proof uses $f(U)$ (the image of an $x^n$-approximator $U$) as the approximator
for the outer product.
\end{theorem}

\begin{theorem}
\label{thm.fps.prods-mulable-fibers-eq.lean}
\lean{PowerSeries.tprod_eq_tprod_fibers}
\leanhelper
\leanok
Under the hypotheses of Theorem~\ref{thm.fps.prods-mulable-fibers.lean},
\[
\prod_{s \in S} \mathbf{a}_s
= \prod_{w \in W} \Bigl(\prod_{s \in f^{-1}(w)} \mathbf{a}_s\Bigr).
\]
This is the Lean formalization of Proposition~\ref{prop.fps.prods-mulable-rules.SW1}
(parts \textbf{(a)} and \textbf{(b)}).
\end{theorem}

\subsection{Fubini rules for infinite products}

\begin{theorem}
\label{thm.fps.fubini-mulable.lean}
\lean{PowerSeries.multipliable_tprod_fubini}
\leanhelper
\leanok
Let $(\mathbf{a}_{(i,j)})_{(i,j) \in I \times J}$ be a multipliable family.
Suppose that for each $i \in I$, the row family $(\mathbf{a}_{(i,j)})_{j \in J}$
is multipliable. Then the family of row products
$\bigl(\prod_{j \in J} \mathbf{a}_{(i,j)}\bigr)_{i \in I}$ is multipliable.
\end{theorem}

\begin{theorem}
\label{thm.fps.fubini.lean}
\lean{PowerSeries.tprod_fubini}
\leanhelper
\leanok
Under the hypotheses of Theorem~\ref{thm.fps.fubini-mulable.lean},
\[
\prod_{(i,j) \in I \times J} \mathbf{a}_{(i,j)}
= \prod_{i \in I} \Bigl(\prod_{j \in J} \mathbf{a}_{(i,j)}\Bigr).
\]
This is the Fubini rule for infinite products of FPSs.
\end{theorem}

\begin{theorem}
\label{thm.fps.fubini-inv.lean}
\lean{PowerSeries.fubini_prod_invertible}
\leanhelper
\leanok
If $(\mathbf{a}_{(i,j)})_{(i,j) \in I \times J}$ is a multipliable family
of invertible FPSs (each has unit constant term), then both the row families
$(\mathbf{a}_{(i,j)})_{j \in J}$ for each $i$ and the column families
$(\mathbf{a}_{(i,j)})_{i \in I}$ for each $j$ are multipliable. No additional
multipliability hypotheses are needed in the invertible case, because
subfamilies of invertible multipliable families are automatically multipliable.
\end{theorem}

\subsection{Coefficient form of the generalized distributive law}

\begin{lemma}
\label{lem.fps.prodrule-inf-inf-coeff}
\lean{PowerSeries.prodRule_infInf_coeff}
\leanhelper
\leanok
For any finite set $M \subseteq I$ and finite value sets $(S_i)_{i \in M}$,
\[
[x^n]\Bigl(\prod_{i \in M} \sum_{k \in S_i} p_{i,k}\Bigr)
= [x^n]\Bigl(\sum_{(k_i) \in \prod_{i \in M} S_i} \prod_{i \in M} p_{i,k_i}\Bigr).
\]
This is a direct corollary of Claim~11 (finite product-sum interchange).
\end{lemma}

\subsection{Congruence lemma for division}

\begin{lemma}
\label{lem.fps.prod.irlv.cong-div}
\lean{AlgebraicCombinatorics.FPS.xnEquiv_div}
\leantarget
\leanok
Let $a, b, c, d \in K\left[\left[x\right]\right]$ be four FPSs such that
$c$ and $d$ are invertible. Let $n \in \mathbb{N}$. Assume that
\[
\left[x^{m}\right] a = \left[x^{m}\right] b
\qquad \text{for each } m \in \left\{0, 1, \ldots, n\right\}.
\]
Assume further that
\[
\left[x^{m}\right] c = \left[x^{m}\right] d
\qquad \text{for each } m \in \left\{0, 1, \ldots, n\right\}.
\]
Then,
\[
\left[x^{m}\right] \dfrac{a}{c} = \left[x^{m}\right] \dfrac{b}{d}
\qquad \text{for each } m \in \left\{0, 1, \ldots, n\right\}.
\]
\end{lemma}

\begin{proof}
\leanok
We have assumed that
\[
\left[x^{m}\right] a = \left[x^{m}\right] b
\qquad \text{for each } m \in \left\{0, 1, \ldots, n\right\}.
\]
In other words, $a \overset{x^n}{\equiv} b$ (by the definition of the
relation $\overset{x^n}{\equiv}$). Similarly,
$c \overset{x^n}{\equiv} d$. Hence, by the compatibility of
$\overset{x^n}{\equiv}$ with division
(Theorem~\ref{thm.fps.xneq.props}~\textbf{(e)}), we obtain
$\dfrac{a}{c} \overset{x^n}{\equiv} \dfrac{b}{d}$. In other words,
\[
\left[x^{m}\right] \dfrac{a}{c} = \left[x^{m}\right] \dfrac{b}{d}
\qquad \text{for each } m \in \left\{0, 1, \ldots, n\right\}
\]
(by the definition of the relation $\overset{x^n}{\equiv}$).
This proves Lemma~\ref{lem.fps.prod.irlv.cong-div}.
\end{proof}

\subsection{Detailed proof of Proposition~\ref{prop.fps.div-mulable}}

\begin{proof}[Detailed proof of Proposition~\ref{prop.fps.div-mulable}]
\proves{prop.fps.div-mulable}
\leanok
This can be proved similarly to Proposition~\ref{prop.fps.prod-mulable},
with the obvious changes (replacing multiplication by division, and
requiring each $\mathbf{b}_i$ to be invertible). Of course, instead of
Lemma~\ref{lem.fps.prod.irlv.cong-mul}, we need to use
Lemma~\ref{lem.fps.prod.irlv.cong-div}.
\end{proof}

\subsection{Approximator intersection lemma}

\begin{lemma}
\label{lem.fps.prods-mulable-subfams-appr}
\lean{AlgebraicCombinatorics.FPS.xnApproximator_inter}
\leantarget
\leanok
Let $\left(\mathbf{a}_{i}\right)_{i \in I} \in
K\left[\left[x\right]\right]^{I}$ be a family of invertible FPSs.
Let $J$ be a subset of $I$. Let $n \in \mathbb{N}$. Let $U$ be an
$x^n$-approximator for $\left(\mathbf{a}_{i}\right)_{i \in I}$.
Then, $U \cap J$ is an $x^n$-approximator for
$\left(\mathbf{a}_{i}\right)_{i \in J}$.
\end{lemma}

\begin{proof}
\leanok
The set $U$ is an $x^n$-approximator for
$\left(\mathbf{a}_{i}\right)_{i \in I}$. In other words, $U$ is a
finite subset of $I$ that determines the first $n+1$ coefficients in the
product of $\left(\mathbf{a}_{i}\right)_{i \in I}$ (by the definition
of an $x^n$-approximator). Hence, in particular, $U$ is finite.
Moreover, $U \subseteq I$ (since $U$ is a subset of $I$).

Let $M = U \cap J$. Thus, $M = U \cap J \subseteq U$, so that the set
$M$ is finite (since $U$ is finite). Moreover, $M$ is a subset of $J$
(since $M = U \cap J \subseteq J$).

Now, let $N$ be a finite subset of $J$ satisfying
$M \subseteq N \subseteq J$. We shall show that
\[
\prod_{i \in N} \mathbf{a}_{i} \overset{x^n}{\equiv}
\prod_{i \in M} \mathbf{a}_{i}.
\]

Indeed, the set $N \cup U$ is finite (since $N$ and $U$ are finite) and
is a subset of $I$ (since $N \subseteq J \subseteq I$ and
$U \subseteq I$); it also satisfies $U \subseteq N \cup U \subseteq I$.
Now, we recall that $U$ is an $x^n$-approximator for
$\left(\mathbf{a}_{i}\right)_{i \in I}$. Hence,
Proposition~\ref{prop.fps.infprod-approx-xneq}~\textbf{(a)} (applied to
$U$ and $N \cup U$ instead of $M$ and $J$) yields
\[
\prod_{i \in N \cup U} \mathbf{a}_{i} \overset{x^n}{\equiv}
\prod_{i \in U} \mathbf{a}_{i}
\]
(since $N \cup U$ is a finite subset of $I$ satisfying
$U \subseteq N \cup U \subseteq I$).

On the other hand, we have assumed that
$\left(\mathbf{a}_{i}\right)_{i \in I}$ is a family of invertible FPSs.
Thus, for each $i \in I$, the FPS $\mathbf{a}_{i}$ is invertible, so
that its inverse $\mathbf{a}_{i}^{-1}$ is well-defined. We obviously
have
\[
\prod_{i \in U \setminus J} \mathbf{a}_{i}^{-1}
\overset{x^n}{\equiv}
\prod_{i \in U \setminus J} \mathbf{a}_{i}^{-1}
\]
(since the relation $\overset{x^n}{\equiv}$ is reflexive).

Hence, by the compatibility of $\overset{x^n}{\equiv}$ with
multiplication
(Theorem~\ref{thm.fps.xneq.props}~\textbf{(b)}), applied to
$a = \prod_{i \in N \cup U} \mathbf{a}_{i}$ and
$b = \prod_{i \in U} \mathbf{a}_{i}$ and
$c = \prod_{i \in U \setminus J} \mathbf{a}_{i}^{-1}$ and
$d = \prod_{i \in U \setminus J} \mathbf{a}_{i}^{-1}$, we obtain
\[
\left(\prod_{i \in N \cup U} \mathbf{a}_{i}\right)
\left(\prod_{i \in U \setminus J} \mathbf{a}_{i}^{-1}\right)
\overset{x^n}{\equiv}
\left(\prod_{i \in U} \mathbf{a}_{i}\right)
\left(\prod_{i \in U \setminus J} \mathbf{a}_{i}^{-1}\right).
\]

On the other hand, the sets $N$ and $U \setminus J$ are disjoint
(since $N \subseteq J$ and $U \setminus J \subseteq J^c$), and their
union is $N \cup (U \setminus J) = N \cup U$
(since $U = (U \cap J) \cup (U \setminus J)$ and
$U \cap J = M \subseteq N$). Hence, the set $N \cup U$ is the union of
its two disjoint subsets $N$ and $U \setminus J$. Thus, we can split the
product $\prod_{i \in N \cup U} \mathbf{a}_{i}$ as follows:
\[
\prod_{i \in N \cup U} \mathbf{a}_{i}
= \left(\prod_{i \in N} \mathbf{a}_{i}\right)
  \left(\prod_{i \in U \setminus J} \mathbf{a}_{i}\right).
\]
Multiplying both sides of this equality by
$\prod_{i \in U \setminus J} \mathbf{a}_{i}^{-1}$, we obtain
\begin{align*}
\left(\prod_{i \in N \cup U} \mathbf{a}_{i}\right)
\left(\prod_{i \in U \setminus J} \mathbf{a}_{i}^{-1}\right)
&= \left(\prod_{i \in N} \mathbf{a}_{i}\right)
   \underbrace{
     \left(\prod_{i \in U \setminus J} \mathbf{a}_{i}\right)
     \left(\prod_{i \in U \setminus J} \mathbf{a}_{i}^{-1}\right)
   }_{= \prod_{i \in U \setminus J}
     \left(\mathbf{a}_{i} \mathbf{a}_{i}^{-1}\right)} \\
&= \left(\prod_{i \in N} \mathbf{a}_{i}\right)
   \left(\prod_{i \in U \setminus J}
     \underbrace{\left(\mathbf{a}_{i} \mathbf{a}_{i}^{-1}\right)}_{= 1}\right)
 = \left(\prod_{i \in N} \mathbf{a}_{i}\right)
   \underbrace{\left(\prod_{i \in U \setminus J} 1\right)}_{= 1} \\
&= \prod_{i \in N} \mathbf{a}_{i}.
\end{align*}

However, the set $U$ is the union of its two disjoint subsets $U \cap J$
and $U \setminus J$. Thus, we can split the product
$\prod_{i \in U} \mathbf{a}_{i}$ as follows:
\[
\prod_{i \in U} \mathbf{a}_{i}
= \left(\prod_{i \in U \cap J} \mathbf{a}_{i}\right)
  \left(\prod_{i \in U \setminus J} \mathbf{a}_{i}\right).
\]
Multiplying both sides by
$\prod_{i \in U \setminus J} \mathbf{a}_{i}^{-1}$, we obtain
\begin{align*}
\left(\prod_{i \in U} \mathbf{a}_{i}\right)
\left(\prod_{i \in U \setminus J} \mathbf{a}_{i}^{-1}\right)
&= \underbrace{\left(\prod_{i \in U \cap J} \mathbf{a}_{i}\right)}_{
     = \prod_{i \in M} \mathbf{a}_{i}
     \text{ (since } U \cap J = M \text{)}}
   \underbrace{
     \left(\prod_{i \in U \setminus J} \mathbf{a}_{i}\right)
     \left(\prod_{i \in U \setminus J} \mathbf{a}_{i}^{-1}\right)
   }_{= \prod_{i \in U \setminus J}
     \left(\mathbf{a}_{i} \mathbf{a}_{i}^{-1}\right)} \\
&= \left(\prod_{i \in M} \mathbf{a}_{i}\right)
   \left(\prod_{i \in U \setminus J}
     \underbrace{\left(\mathbf{a}_{i} \mathbf{a}_{i}^{-1}\right)}_{= 1}\right)
 = \left(\prod_{i \in M} \mathbf{a}_{i}\right)
   \underbrace{\left(\prod_{i \in U \setminus J} 1\right)}_{= 1} \\
&= \prod_{i \in M} \mathbf{a}_{i}.
\end{align*}

Substituting these two simplifications into the
$x^n$-equivalence above, we obtain
\[
\prod_{i \in N} \mathbf{a}_{i} \overset{x^n}{\equiv}
\prod_{i \in M} \mathbf{a}_{i}.
\]
In other words, each $m \in \left\{0, 1, \ldots, n\right\}$ satisfies
\[
\left[x^{m}\right] \left(\prod_{i \in N} \mathbf{a}_{i}\right)
= \left[x^{m}\right] \left(\prod_{i \in M} \mathbf{a}_{i}\right)
\]
(by the definition of $\overset{x^n}{\equiv}$).

Forget that we fixed $N$. We have thus shown that every finite subset
$N$ of $J$ satisfying $M \subseteq N \subseteq J$ satisfies the above
equality for each $m \in \left\{0, 1, \ldots, n\right\}$.

Now, let $m \in \left\{0, 1, \ldots, n\right\}$. Then, every finite
subset $N$ of $J$ satisfying $M \subseteq N \subseteq J$ satisfies
\[
\left[x^{m}\right] \left(\prod_{i \in N} \mathbf{a}_{i}\right)
= \left[x^{m}\right] \left(\prod_{i \in M} \mathbf{a}_{i}\right).
\]
In other words, the set $M$ determines the $x^m$-coefficient in the
product of $\left(\mathbf{a}_{i}\right)_{i \in J}$ (by the definition
of ``determining the $x^m$-coefficient in a product'').

Forget that we fixed $m$. We thus have shown that $M$ determines the
$x^m$-coefficient in the product of
$\left(\mathbf{a}_{i}\right)_{i \in J}$ for each
$m \in \left\{0, 1, \ldots, n\right\}$. In other words, $M$ determines
the first $n+1$ coefficients in the product of
$\left(\mathbf{a}_{i}\right)_{i \in J}$. In other words, $M$ is an
$x^n$-approximator for $\left(\mathbf{a}_{i}\right)_{i \in J}$ (by the
definition of an ``$x^n$-approximator'', since $M$ is a finite subset
of $J$). In other words, $U \cap J$ is an $x^n$-approximator for
$\left(\mathbf{a}_{i}\right)_{i \in J}$ (since $M = U \cap J$).
This proves Lemma~\ref{lem.fps.prods-mulable-subfams-appr}.
\end{proof}

\subsection{Detailed proof of Proposition~\ref{prop.fps.prods-mulable-subfams}}

\begin{proof}[Detailed proof of Proposition~\ref{prop.fps.prods-mulable-subfams}]
\proves{prop.fps.prods-mulable-subfams}
\leanok
Let $J$ be a subset of $I$. We shall show that the family
$\left(\mathbf{a}_{i}\right)_{i \in J}$ is multipliable.

Fix $n \in \mathbb{N}$. We know that the family
$\left(\mathbf{a}_{i}\right)_{i \in I}$ is multipliable. Hence, there
exists an $x^n$-approximator $U$ for
$\left(\mathbf{a}_{i}\right)_{i \in I}$ (by
Lemma~\ref{lem.fps.mulable.approx}). Consider this $U$.

Set $M = U \cap J$. Lemma~\ref{lem.fps.prods-mulable-subfams-appr}
yields that $U \cap J$ is an $x^n$-approximator for
$\left(\mathbf{a}_{i}\right)_{i \in J}$. In other words, $M$ is an
$x^n$-approximator for $\left(\mathbf{a}_{i}\right)_{i \in J}$ (since
$M = U \cap J$). In other words, $M$ is a finite subset of $J$ that
determines the first $n+1$ coefficients in the product of
$\left(\mathbf{a}_{i}\right)_{i \in J}$ (by the definition of an
$x^n$-approximator). Thus, the set $M$ determines the first $n+1$
coefficients in the product of
$\left(\mathbf{a}_{i}\right)_{i \in J}$. Hence, in particular, this set
$M$ determines the $x^n$-coefficient in the product of
$\left(\mathbf{a}_{i}\right)_{i \in J}$. Therefore, the
$x^n$-coefficient in the product of
$\left(\mathbf{a}_{i}\right)_{i \in J}$ is finitely determined (by the
definition of ``finitely determined'', since $M$ is a finite subset
of $J$).

Forget that we fixed $n$. We thus have proved that for each
$n \in \mathbb{N}$, the $x^n$-coefficient in the product of
$\left(\mathbf{a}_{i}\right)_{i \in J}$ is finitely determined. In
other words, each coefficient in the product of
$\left(\mathbf{a}_{i}\right)_{i \in J}$ is finitely determined. In
other words, the family $\left(\mathbf{a}_{i}\right)_{i \in J}$ is
multipliable (by the definition of ``multipliable'').

Forget that we fixed $J$. We thus have shown that the family
$\left(\mathbf{a}_{i}\right)_{i \in J}$ is multipliable whenever $J$ is
a subset of $I$. In other words, any subfamily of
$\left(\mathbf{a}_{i}\right)_{i \in I}$ is multipliable.
This proves Proposition~\ref{prop.fps.prods-mulable-subfams}.
\end{proof}

\subsection{Finite products and $x^n$-equivalence}

\begin{lemma}
\label{lem.fps.prods-mulable-rules.SW1.lem1}
\lean{AlgebraicCombinatorics.FPS.xnEquiv_finprod}
\leantarget
\leanok
Let $n \in \mathbb{N}$. Let $V$ be a finite set. Let
$\left(c_{w}\right)_{w \in V} \in K\left[\left[x\right]\right]^{V}$ and
$\left(d_{w}\right)_{w \in V} \in K\left[\left[x\right]\right]^{V}$ be
two families of FPSs such that
\[
\text{each } w \in V \text{ satisfies }
c_{w} \overset{x^n}{\equiv} d_{w}.
\]
Then, we have
\[
\prod_{w \in V} c_{w} \overset{x^n}{\equiv}
\prod_{w \in V} d_{w}.
\]
\end{lemma}

\begin{proof}
\leanok
This is just Theorem~\ref{thm.fps.xneq.props}~\textbf{(f)} (the
part about products), with the letters $S$, $s$, $a_s$, and $b_s$
renamed as $V$, $w$, $c_w$, and $d_w$.
\end{proof}

\subsection{Detailed proof of Proposition~\ref{prop.fps.prods-mulable-rules.SW1}}

\begin{proof}[Detailed proof of Proposition~\ref{prop.fps.prods-mulable-rules.SW1}]
\proves{prop.fps.prods-mulable-rules.SW1}
\leanok
We shall subdivide our proof into several claims.

\emph{Claim 1:} Let $w \in W$. Then, the family
$\left(\mathbf{a}_{s}\right)_{s \in S,\; f(s) = w}$ is multipliable.

[\emph{Proof of Claim 1:} This was an assumption of
Proposition~\ref{prop.fps.prods-mulable-rules.SW1}.]

Let us set
\[
\mathbf{b}_{w} := \prod_{\substack{s \in S;\\ f(s) = w}}
\mathbf{a}_{s}
\qquad \text{for each } w \in W.
\]
This is well-defined, because for each $w \in W$, the product
$\prod_{\substack{s \in S;\\ f(s) = w}} \mathbf{a}_{s}$ is well-defined
(since Claim~1 shows that the family
$\left(\mathbf{a}_{s}\right)_{s \in S,\; f(s) = w}$ is multipliable).

Now, let $n \in \mathbb{N}$. By Lemma~\ref{lem.fps.mulable.approx}
(applied to $S$ and $\left(\mathbf{a}_{s}\right)_{s \in S}$ instead of
$I$ and $\left(\mathbf{a}_{i}\right)_{i \in I}$), there exists an
$x^n$-approximator for $\left(\mathbf{a}_{s}\right)_{s \in S}$. Pick
such an $x^n$-approximator, and call it $U$. Then, $U$ is an
$x^n$-approximator for $\left(\mathbf{a}_{s}\right)_{s \in S}$; in
other words, $U$ is a finite subset of $S$ that determines the first
$n+1$ coefficients in the product of
$\left(\mathbf{a}_{s}\right)_{s \in S}$ (by the definition of an
$x^n$-approximator).

The set $U$ is finite. Thus, its image
$f(U) = \left\{f(u) \mid u \in U\right\}$ is finite as well. Now, we
claim the following:

\emph{Claim 2:} For each $w \in W$, we have
\[
\mathbf{b}_{w} \overset{x^n}{\equiv}
\prod_{\substack{s \in U;\\ f(s) = w}} \mathbf{a}_{s}.
\]

[\emph{Proof of Claim 2:} Let $w \in W$. Let $J$ be the subset
$\left\{s \in S \mid f(s) = w\right\}$ of $S$. Then, the family
$\left(\mathbf{a}_{s}\right)_{s \in J}$ is just the family
$\left(\mathbf{a}_{s}\right)_{s \in S,\; f(s) = w}$, and thus is
multipliable (by Claim~1). Furthermore,
Lemma~\ref{lem.fps.prods-mulable-subfams-appr} (applied to $S$ and
$\left(\mathbf{a}_{s}\right)_{s \in S}$ instead of $I$ and
$\left(\mathbf{a}_{i}\right)_{i \in I}$) yields that $U \cap J$ is an
$x^n$-approximator for $\left(\mathbf{a}_{s}\right)_{s \in J}$ (since
$U$ is an $x^n$-approximator for
$\left(\mathbf{a}_{s}\right)_{s \in S}$). Hence,
Proposition~\ref{prop.fps.infprod-approx-xneq}~\textbf{(b)} (applied to
$J$, $\left(\mathbf{a}_{s}\right)_{s \in J}$ and $U \cap J$ instead of
$I$, $\left(\mathbf{a}_{i}\right)_{i \in I}$ and $M$) yields
\[
\prod_{s \in J} \mathbf{a}_{s} \overset{x^n}{\equiv}
\prod_{s \in U \cap J} \mathbf{a}_{s}.
\]
However, we have $J = \left\{s \in S \mid f(s) = w\right\}$. Thus,
$\prod_{s \in J} \mathbf{a}_{s}
= \prod_{\substack{s \in S;\\ f(s) = w}} \mathbf{a}_{s} = \mathbf{b}_w$
and $U \cap J = \left\{s \in U \mid f(s) = w\right\}$. So the product
sign ``$\prod_{s \in U \cap J}$'' is equivalent to
``$\prod_{\substack{s \in U;\\ f(s) = w}}$''. Thus,
\[
\mathbf{b}_{w} \overset{x^n}{\equiv}
\prod_{\substack{s \in U;\\ f(s) = w}} \mathbf{a}_{s}.
\]
This proves Claim~2.]

\emph{Claim 3:} The set $f(U)$ determines the $x^n$-coefficient in the
product of $\left(\mathbf{b}_{w}\right)_{w \in W}$.

[\emph{Proof of Claim 3:} Let $T$ be a finite subset of $W$ satisfying
$f(U) \subseteq T$. We must show that
$\left[x^n\right] \left(\prod_{w \in T} \mathbf{b}_w\right)
= \left[x^n\right] \left(\prod_{w \in f(U)} \mathbf{b}_w\right)$.

By Claim~2, for each $w \in W$ we have
$\mathbf{b}_{w} \overset{x^n}{\equiv}
\prod_{\substack{s \in U;\\ f(s) = w}} \mathbf{a}_{s}$.
Hence, by Lemma~\ref{lem.fps.prods-mulable-rules.SW1.lem1}, we get
\[
\prod_{w \in T} \mathbf{b}_w \overset{x^n}{\equiv}
\prod_{w \in T} \prod_{\substack{s \in U;\\ f(s) = w}} \mathbf{a}_{s}
\]
and similarly
\[
\prod_{w \in f(U)} \mathbf{b}_w \overset{x^n}{\equiv}
\prod_{w \in f(U)} \prod_{\substack{s \in U;\\ f(s) = w}} \mathbf{a}_{s}.
\]
Now, for $w \notin f(U)$, there are no $s \in U$ with $f(s) = w$, so
the inner product is $1$ (the empty product). Hence, for any finite set
$V \supseteq f(U)$,
\[
\prod_{w \in V} \prod_{\substack{s \in U;\\ f(s) = w}} \mathbf{a}_{s}
= \prod_{s \in U} \mathbf{a}_{s}
\]
(by the standard fiberwise reindexing of finite products). Applying this
to both $V = T$ and $V = f(U)$, we conclude that both products equal
$\prod_{s \in U} \mathbf{a}_{s}$. Therefore,
\[
\left[x^n\right] \left(\prod_{w \in T} \mathbf{b}_w\right)
= \left[x^n\right] \left(\prod_{s \in U} \mathbf{a}_s\right)
= \left[x^n\right] \left(\prod_{w \in f(U)} \mathbf{b}_w\right).
\]
This proves Claim~3.]

From Claim~3, the $x^n$-coefficient in the product of
$\left(\mathbf{b}_w\right)_{w \in W}$ is finitely determined (since
$f(U)$ is a finite subset of $W$). Since this holds for each
$n \in \mathbb{N}$, the family
$\left(\mathbf{b}_w\right)_{w \in W}$ is multipliable. This proves
part~\textbf{(a)} of Proposition~\ref{prop.fps.prods-mulable-rules.SW1}.

For part~\textbf{(b)}, we observe that by
Proposition~\ref{prop.fps.infprod-approx-xneq}~\textbf{(b)}, the
infinite product $\prod_{s \in S} \mathbf{a}_s$ is $x^n$-equivalent to
$\prod_{s \in U} \mathbf{a}_s$. By the computation above, the infinite
product $\prod_{w \in W} \mathbf{b}_w$ is also $x^n$-equivalent to
$\prod_{s \in U} \mathbf{a}_s$. Since this holds for all $n$, the two
infinite products are equal.
\end{proof}

\subsection{Detailed proof of Proposition~\ref{prop.fps.prod-mulable}}

\begin{proof}[Detailed proof of Proposition~\ref{prop.fps.prod-mulable}]
\proves{prop.fps.prod-mulable}
\leanok
\textbf{(a)} Fix $n \in \mathbb{N}$. We know that the family
$\left(\mathbf{a}_{i}\right)_{i \in I}$ is multipliable. Hence, there
exists an $x^n$-approximator $U$ for
$\left(\mathbf{a}_{i}\right)_{i \in I}$ (by
Lemma~\ref{lem.fps.mulable.approx}). Consider this $U$.

We also know that the family $\left(\mathbf{b}_{i}\right)_{i \in I}$ is
multipliable. Hence, there exists an $x^n$-approximator $V$ for
$\left(\mathbf{b}_{i}\right)_{i \in I}$ (by
Lemma~\ref{lem.fps.mulable.approx}, applied to $\mathbf{b}_{i}$ instead
of $\mathbf{a}_{i}$). Consider this $V$.

We know that $U$ is an $x^n$-approximator for
$\left(\mathbf{a}_{i}\right)_{i \in I}$. In other words, $U$ is a
finite subset of $I$ that determines the first $n+1$ coefficients in the
product of $\left(\mathbf{a}_{i}\right)_{i \in I}$ (by the definition
of an $x^n$-approximator). Hence, in particular, $U$ is finite.
Similarly, $V$ is finite. Moreover, $U \subseteq I$; similarly,
$V \subseteq I$.

Let $M = U \cup V$. This set $M$ is finite (since $U$ and $V$ are
finite). Moreover, $M = U \cup V \subseteq I \cup I = I$. Hence, $M$ is
a finite subset of $I$.

Now, let $N$ be a finite subset of $I$ satisfying
$M \subseteq N \subseteq I$. We shall show that
\[
\left[x^{n}\right] \left(\prod_{i \in N}
  \left(\mathbf{a}_{i} \mathbf{b}_{i}\right)\right)
= \left[x^{n}\right] \left(\prod_{i \in M}
  \left(\mathbf{a}_{i} \mathbf{b}_{i}\right)\right).
\]

Indeed, let $m \in \left\{0, 1, \ldots, n\right\}$. We have
$U \subseteq U \cup V = M \subseteq N$. The set $N$ is a finite subset
of $I$ and satisfies $U \subseteq N$. Now, recall that $U$ determines
the first $n+1$ coefficients in the product of
$\left(\mathbf{a}_{i}\right)_{i \in I}$. Hence, $U$ determines the
$x^m$-coefficient in the product of
$\left(\mathbf{a}_{i}\right)_{i \in I}$ (since
$m \in \left\{0, 1, \ldots, n\right\}$). In other words, every finite
subset $T$ of $I$ satisfying $U \subseteq T \subseteq I$ satisfies
\[
\left[x^{m}\right] \left(\prod_{i \in T} \mathbf{a}_{i}\right)
= \left[x^{m}\right] \left(\prod_{i \in U} \mathbf{a}_{i}\right).
\]
We can apply this to $T = N$ (since $N$ is a finite subset of $I$
satisfying $U \subseteq N \subseteq I$), and thus obtain
\[
\left[x^{m}\right] \left(\prod_{i \in N} \mathbf{a}_{i}\right)
= \left[x^{m}\right] \left(\prod_{i \in U} \mathbf{a}_{i}\right).
\]
However, we can also apply it to $T = M$ (since $M$ is a finite subset
of $I$ satisfying $U \subseteq M \subseteq I$), and thus obtain
\[
\left[x^{m}\right] \left(\prod_{i \in M} \mathbf{a}_{i}\right)
= \left[x^{m}\right] \left(\prod_{i \in U} \mathbf{a}_{i}\right).
\]
Comparing these two equalities, we obtain
\[
\left[x^{m}\right] \left(\prod_{i \in N} \mathbf{a}_{i}\right)
= \left[x^{m}\right] \left(\prod_{i \in M} \mathbf{a}_{i}\right).
\]
The same argument (applied to $\mathbf{b}_{i}$ and $V$ instead of
$\mathbf{a}_{i}$ and $U$) yields
\[
\left[x^{m}\right] \left(\prod_{i \in N} \mathbf{b}_{i}\right)
= \left[x^{m}\right] \left(\prod_{i \in M} \mathbf{b}_{i}\right).
\]

Forget that we fixed $m$. We thus have proved these equalities for each
$m \in \left\{0, 1, \ldots, n\right\}$. Hence,
Lemma~\ref{lem.fps.prod.irlv.cong-mul} (applied to
$a = \prod_{i \in N} \mathbf{a}_{i}$ and
$b = \prod_{i \in M} \mathbf{a}_{i}$ and
$c = \prod_{i \in N} \mathbf{b}_{i}$ and
$d = \prod_{i \in M} \mathbf{b}_{i}$) yields that
\[
\left[x^{m}\right] \left(
  \left(\prod_{i \in N} \mathbf{a}_{i}\right)
  \left(\prod_{i \in N} \mathbf{b}_{i}\right)\right)
= \left[x^{m}\right] \left(
  \left(\prod_{i \in M} \mathbf{a}_{i}\right)
  \left(\prod_{i \in M} \mathbf{b}_{i}\right)\right)
\]
for each $m \in \left\{0, 1, \ldots, n\right\}$. Applying this to
$m = n$, we obtain
\[
\left[x^{n}\right] \left(
  \left(\prod_{i \in N} \mathbf{a}_{i}\right)
  \left(\prod_{i \in N} \mathbf{b}_{i}\right)\right)
= \left[x^{n}\right] \left(
  \left(\prod_{i \in M} \mathbf{a}_{i}\right)
  \left(\prod_{i \in M} \mathbf{b}_{i}\right)\right).
\]
In view of
\[
\prod_{i \in N} \left(\mathbf{a}_{i} \mathbf{b}_{i}\right)
= \left(\prod_{i \in N} \mathbf{a}_{i}\right)
  \left(\prod_{i \in N} \mathbf{b}_{i}\right)
\qquad \text{and} \qquad
\prod_{i \in M} \left(\mathbf{a}_{i} \mathbf{b}_{i}\right)
= \left(\prod_{i \in M} \mathbf{a}_{i}\right)
  \left(\prod_{i \in M} \mathbf{b}_{i}\right),
\]
this rewrites as
\[
\left[x^{n}\right] \left(\prod_{i \in N}
  \left(\mathbf{a}_{i} \mathbf{b}_{i}\right)\right)
= \left[x^{n}\right] \left(\prod_{i \in M}
  \left(\mathbf{a}_{i} \mathbf{b}_{i}\right)\right).
\]

Forget that we fixed $N$. We thus have shown that every finite subset
$N$ of $I$ satisfying $M \subseteq N \subseteq I$ satisfies the above
equality. In other words, $M$ determines the $x^n$-coefficient in the
product of $\left(\mathbf{a}_{i} \mathbf{b}_{i}\right)_{i \in I}$ (by
the definition of ``determining the $x^n$-coefficient''). Hence, the
$x^n$-coefficient in the product of
$\left(\mathbf{a}_{i} \mathbf{b}_{i}\right)_{i \in I}$ is finitely
determined (since $M$ is a finite subset of $I$).

Forget that we fixed $n$. We thus have proved that for each
$n \in \mathbb{N}$, the $x^n$-coefficient in the product of
$\left(\mathbf{a}_{i} \mathbf{b}_{i}\right)_{i \in I}$ is finitely
determined. In other words, the family
$\left(\mathbf{a}_{i} \mathbf{b}_{i}\right)_{i \in I}$ is multipliable.
This proves Proposition~\ref{prop.fps.prod-mulable}~\textbf{(a)}.

\medskip

\textbf{(b)} Proposition~\ref{prop.fps.prod-mulable}~\textbf{(a)} shows
that the family
$\left(\mathbf{a}_{i} \mathbf{b}_{i}\right)_{i \in I}$ is multipliable.

Let $n \in \mathbb{N}$. In our above proof of
Proposition~\ref{prop.fps.prod-mulable}~\textbf{(a)}, we have seen the
following:
\begin{itemize}
\item There exists an $x^n$-approximator $U$ for
  $\left(\mathbf{a}_{i}\right)_{i \in I}$.
\item There exists an $x^n$-approximator $V$ for
  $\left(\mathbf{b}_{i}\right)_{i \in I}$.
\end{itemize}
Consider these $U$ and $V$. Let $M = U \cup V$. Thus,
$M = U \cup V \supseteq U$ and $M = U \cup V \supseteq V$. In our above
proof, we have seen that $M$ is a finite subset of $I$ and $M$
determines the $x^n$-coefficient in the product of
$\left(\mathbf{a}_{i} \mathbf{b}_{i}\right)_{i \in I}$.

We recall that the relation $\overset{x^n}{\equiv}$ is an equivalence
relation (by Theorem~\ref{thm.fps.xneq.props}~\textbf{(a)}).

Now, the definition of the infinite product
$\prod_{i \in I} \left(\mathbf{a}_{i} \mathbf{b}_{i}\right)$ (namely,
Definition~\ref{def.fps.multipliable}~\textbf{(b)}) yields that
\[
\left[x^{n}\right] \left(\prod_{i \in I}
  \left(\mathbf{a}_{i} \mathbf{b}_{i}\right)\right)
= \left[x^{n}\right] \left(\prod_{i \in M}
  \left(\mathbf{a}_{i} \mathbf{b}_{i}\right)\right)
\]
(since $M$ is a finite subset of $I$ that determines the
$x^n$-coefficient in the product of
$\left(\mathbf{a}_{i} \mathbf{b}_{i}\right)_{i \in I}$). On the other
hand, the set $U$ is an $x^n$-approximator for
$\left(\mathbf{a}_{i}\right)_{i \in I}$. Thus,
Proposition~\ref{prop.fps.infprod-approx-xneq}~\textbf{(b)} (applied to
$U$ instead of $M$) yields
\[
\prod_{i \in I} \mathbf{a}_{i} \overset{x^n}{\equiv}
\prod_{i \in U} \mathbf{a}_{i}.
\]
Since the relation $\overset{x^n}{\equiv}$ is symmetric, we thus obtain
$\prod_{i \in U} \mathbf{a}_{i} \overset{x^n}{\equiv}
\prod_{i \in I} \mathbf{a}_{i}$.
Moreover, $M$ is a finite subset of $I$ satisfying $U \subseteq M$;
hence, Proposition~\ref{prop.fps.infprod-approx-xneq}~\textbf{(a)}
yields
\[
\prod_{i \in M} \mathbf{a}_{i} \overset{x^n}{\equiv}
\prod_{i \in U} \mathbf{a}_{i} \overset{x^n}{\equiv}
\prod_{i \in I} \mathbf{a}_{i}.
\]
Hence,
$\prod_{i \in M} \mathbf{a}_{i} \overset{x^n}{\equiv}
\prod_{i \in I} \mathbf{a}_{i}$
(since the relation $\overset{x^n}{\equiv}$ is transitive). The same
argument (applied to $\left(\mathbf{b}_{i}\right)_{i \in I}$ and $V$
instead of $\left(\mathbf{a}_{i}\right)_{i \in I}$ and $U$) yields
$\prod_{i \in M} \mathbf{b}_{i} \overset{x^n}{\equiv}
\prod_{i \in I} \mathbf{b}_{i}$.
From these two $x^n$-equivalences, we obtain
\[
\left(\prod_{i \in M} \mathbf{a}_{i}\right)
\left(\prod_{i \in M} \mathbf{b}_{i}\right)
\overset{x^n}{\equiv}
\left(\prod_{i \in I} \mathbf{a}_{i}\right)
\left(\prod_{i \in I} \mathbf{b}_{i}\right)
\]
(by Theorem~\ref{thm.fps.xneq.props}~\textbf{(b)}). In view of
\[
\left(\prod_{i \in M} \mathbf{a}_{i}\right)
\left(\prod_{i \in M} \mathbf{b}_{i}\right)
= \prod_{i \in M}
  \left(\mathbf{a}_{i} \mathbf{b}_{i}\right),
\]
this rewrites as
\[
\prod_{i \in M} \left(\mathbf{a}_{i} \mathbf{b}_{i}\right)
\overset{x^n}{\equiv}
\left(\prod_{i \in I} \mathbf{a}_{i}\right)
\left(\prod_{i \in I} \mathbf{b}_{i}\right).
\]
In other words,
\[
\text{each } m \in \left\{0, 1, \ldots, n\right\} \text{ satisfies }
\left[x^{m}\right] \left(\prod_{i \in M}
  \left(\mathbf{a}_{i} \mathbf{b}_{i}\right)\right)
= \left[x^{m}\right] \left(
  \left(\prod_{i \in I} \mathbf{a}_{i}\right)
  \left(\prod_{i \in I} \mathbf{b}_{i}\right)\right).
\]
Applying this to $m = n$, we obtain
\[
\left[x^{n}\right] \left(\prod_{i \in M}
  \left(\mathbf{a}_{i} \mathbf{b}_{i}\right)\right)
= \left[x^{n}\right] \left(
  \left(\prod_{i \in I} \mathbf{a}_{i}\right)
  \left(\prod_{i \in I} \mathbf{b}_{i}\right)\right).
\]
In view of
$\left[x^{n}\right] \left(\prod_{i \in I}
  \left(\mathbf{a}_{i} \mathbf{b}_{i}\right)\right)
= \left[x^{n}\right] \left(\prod_{i \in M}
  \left(\mathbf{a}_{i} \mathbf{b}_{i}\right)\right)$,
this rewrites as
\[
\left[x^{n}\right] \left(\prod_{i \in I}
  \left(\mathbf{a}_{i} \mathbf{b}_{i}\right)\right)
= \left[x^{n}\right] \left(
  \left(\prod_{i \in I} \mathbf{a}_{i}\right)
  \left(\prod_{i \in I} \mathbf{b}_{i}\right)\right).
\]

Now, forget that we fixed $n$. We thus have proved that each
$n \in \mathbb{N}$ satisfies the above equality. In other words, each
coefficient of the FPS
$\prod_{i \in I} \left(\mathbf{a}_{i} \mathbf{b}_{i}\right)$ equals the
corresponding coefficient of
$\left(\prod_{i \in I} \mathbf{a}_{i}\right)
\left(\prod_{i \in I} \mathbf{b}_{i}\right)$.
Hence, we have
\[
\prod_{i \in I} \left(\mathbf{a}_{i} \mathbf{b}_{i}\right)
= \left(\prod_{i \in I} \mathbf{a}_{i}\right)
  \cdot \left(\prod_{i \in I} \mathbf{b}_{i}\right).
\]
This proves Proposition~\ref{prop.fps.prod-mulable}~\textbf{(b)}.
\end{proof}

\subsection{Product rule claims}

\begin{lemma}
\label{lem.fps.infprod.claim3}
\lean{PowerSeries.prodRule_claim3}
\leanhelper
\leanok
\textbf{(Claim 3.)}
If $(i,k) \notin T'_n$ (the coefficient support union up to $n$) and $k \neq 0$,
then $[x^m] p_{i,k} = 0$ for all $m \in \{0,1,\ldots,n\}$.
\end{lemma}

\begin{proof}
\leanok
If $[x^m] p_{i,k} \neq 0$ for some $m \leq n$, then $(i,k) \in T_m \subseteq T'_n$,
contradicting the hypothesis.
\end{proof}

\begin{lemma}
\label{lem.fps.infprod.claim4}
\lean{PowerSeries.prodRule_claim4}
\leanhelper
\leanok
\textbf{(Claim 4.)}
If $(k_i)_{i \in I}$ is essentially finite and $p_{i,0} = 1$ for all $i$,
then the family $(p_{i,k_i})_{i \in I}$ is multipliable.
\end{lemma}

\begin{proof}
\leanok
All but finitely many $i \in I$ satisfy $k_i = 0$ (since $(k_i)_{i \in I}$
is essentially finite) and thus $p_{i,k_i} = p_{i,0} = 1$.
Hence, all but finitely many entries of the family $(p_{i,k_i})_{i \in I}$
equal $1$. Thus, this family is multipliable (by
Proposition~\ref{prop.fps.1-mulable}).
\end{proof}

\begin{lemma}
\label{lem.fps.infprod.claim5}
\lean{PowerSeries.prodRule_claim5}
\leanhelper
\leanok
\textbf{(Claim 5.)}
If some $j \in I$ satisfies $k_j \neq 0$ and $(j, k_j) \notin T'_n$, then
$[x^n](\prod_{i \in s} p_{i,k_i}) = 0$ for any finite $s \ni j$.
\end{lemma}

\begin{proof}
\leanok
The product $\prod_{i \in s} p_{i,k_i}$ is a multiple of $p_{j,k_j}$
(since $p_{j,k_j}$ is one of the factors of this product).
Moreover, by Claim~3, we have $[x^m] p_{j,k_j} = 0$ for each
$m \in \{0,1,\ldots,n\}$.
Hence, Lemma~\ref{lem.fps.prod.irlv.mul-dvd} (applied to $u = p_{j,k_j}$
and $v = \prod_{i \in s} p_{i,k_i}$) shows that
$[x^m](\prod_{i \in s} p_{i,k_i}) = 0$ for each $m \in \{0,1,\ldots,n\}$.
Applying this to $m = n$ yields the claim.
\end{proof}

\begin{lemma}
\label{lem.fps.infprod.claim6}
\lean{PowerSeries.prodRule_claim6}
\leanhelper
\leanok
\textbf{(Claim 6.)}
If $j \notin I_n$ and $k_j \neq 0$, then
$[x^n](\prod_{i \in s} p_{i,k_i}) = 0$ for any finite $s \ni j$.
\end{lemma}

\begin{proof}
\leanok
We have $k_j \in S_j$ and $k_j \neq 0$, so $(j, k_j) \in \overline{S}$.
Moreover, $(j, k_j) \notin T'_n$ (because if we had
$(j, k_j) \in T'_n$, then we would have $j \in I_n$ by the
definition of $I_n$; but this would contradict $j \notin I_n$).
Hence, $(j, k_j) \in \overline{S} \setminus T'_n$.
Therefore, Claim~5 yields $[x^n](\prod_{i \in s} p_{i,k_i}) = 0$.
\end{proof}

\begin{lemma}
\label{lem.fps.infprod.claim7}
\lean{PowerSeries.prodRule_claim7}
\leanhelper
\leanok
\textbf{(Claim 7.)}
The family $(\prod_{i \in I_n} p_{i,k_i})$ indexed by essentially finite
families $(k_i)$ supported on $I_n$ is summable: for each $n$, only finitely
many such families yield $[x^n](\prod_{i \in I_n} p_{i,k_i}) \neq 0$.
\end{lemma}

\begin{proof}
\leanok
Let $(k_i)_{i \in I} \in S^I_{\mathrm{fin}}$ satisfy
$[x^n](\prod_{i \in I} p_{i,k_i}) \neq 0$.
Then it must satisfy two properties:
(1)~all $i \in I \setminus I_n$ satisfy $k_i = 0$
(otherwise Claim~6 would give a contradiction);
(2)~all $i \in I_n$ satisfy $k_i \in K_n \cup \{0\}$
(otherwise $(j, k_j) \in \overline{S} \setminus T'_n$ for some
$j \in I_n$, and Claim~5 would give a contradiction).
These two properties together restrict $(k_i)$ to finitely many
possibilities (since the sets $I_n$ and $K_n \cup \{0\}$ are finite).
Hence, all but finitely many families yield
$[x^n](\prod_{i \in I} p_{i,k_i}) = 0$, proving summability.
\end{proof}

\begin{lemma}
\label{lem.fps.infprod.claim9}
\lean{PowerSeries.prodRule_claim9}
\leanhelper
\leanok
\textbf{(Claim 9.)}
For $i \notin I_n$, we have $[x^m](\sum_{k \in S_i \setminus \{0\}} p_{i,k}) = 0$
for each $m \in \{0,1,\ldots,n\}$.
\end{lemma}

\begin{proof}
\leanok
Each term $p_{i,k}$ with $k \neq 0$ has $[x^m] p_{i,k} = 0$ by Claim~3
(since $(i,k) \notin T'_n$ because $i \notin I_n$).
\end{proof}

\begin{lemma}
\label{lem.fps.infprod.claim10}
\lean{PowerSeries.prodRule_claim10}
\leanhelper
\leanok
\textbf{(Claim 10.)}
For any finite set $J \supseteq I_n$,
$[x^n](\prod_{i \in J} \sum_{k \in S_i} p_{i,k}) = [x^n](\prod_{i \in I_n} \sum_{k \in S_i} p_{i,k})$.
The product stabilizes at $I_n$.
\end{lemma}

\begin{proof}
\leanok
For $i \in J \setminus I_n$, we have $\sum_{k \in S_i} p_{i,k} = 1 + \sum_{k \in S_i \setminus \{0\}} p_{i,k}$.
By Claim~9, the remainder has zero coefficients up to degree $n$, so multiplying by it
does not change the coefficient at degree $n$.
\end{proof}

\begin{lemma}
\label{lem.fps.infprod.claim11}
\lean{PowerSeries.prodRule_claim11}
\leanhelper
\leanok
\textbf{(Claim 11.)} For a finite index set $I_n$,
\[
\prod_{i \in I_n} \sum_{k \in S_i} p_{i,k} =
\sum_{(k_i) \in S^{I_n}} \prod_{i \in I_n} p_{i,k_i}.
\]
This is the standard finite distributive law (product of sums equals sum of products).
\end{lemma}

\begin{proof}
\leanok
This is the standard finite distributive law
(product of sums equals sum of products over all choice functions).
See Proposition~\ref{prop.fps.prodrule-fin-infJ}.
\end{proof}

\begin{lemma}
\label{lem.fps.infprod.claim8}
\lean{PowerSeries.prodRule_claim8}
\leanhelper
\leanok
\textbf{(Claim 8.)}
The coefficient of the sum over essentially finite families equals the coefficient
of the sum restricted to the finite index set $I_n$:
\[
[x^n]\Bigl(\sum_{\substack{(k_i) \in S^I_{\mathrm{fin}}}} \prod_{i \in I_n} p_{i,k_i}\Bigr)
= [x^n]\Bigl(\sum_{(k_i) \in S^{I_n}} \prod_{i \in I_n} p_{i,k_i}\Bigr).
\]
\end{lemma}

\begin{proof}
\leanok
The proof establishes a bijection between essentially finite families supported on $I_n$
and functions $I_n \to S$, via restriction and extension by $0$.
\end{proof}

\subsection{Composition rules for infinite products}

\begin{lemma}
\label{lem.fps.infprod.xnEquiv-comp}
\lean{PowerSeries.xnEquiv_comp}
\leanhelper
\leanok
If $f_1 \overset{x^n}{\equiv} f_2$ and $[x^0]g = 0$, then
$f_1 \circ g \overset{x^n}{\equiv} f_2 \circ g$.
That is, $x^n$-equivalence is preserved under composition with a fixed inner series.
\end{lemma}

\begin{proof}
\leanok
The coefficient $[x^k](f \circ g)$ involves a sum over $d$ of $[x^d]f \cdot [x^k](g^d)$.
For $d > k$, $[x^k](g^d) = 0$ since $\operatorname{ord}(g^d) \geq d > k$.
For $d \leq k \leq n$, $[x^d]f_1 = [x^d]f_2$ by $x^n$-equivalence.
\end{proof}

\begin{lemma}
\label{lem.fps.infprod.comp-approx}
\lean{PowerSeries.comp_prod_approx_determines}
\leanhelper
\leanok
If $M$ is an $x^n$-approximator for $(f_i)_{i \in I}$ (i.e., for all $k \leq n$
and $J \supseteq M$, $[x^k](\prod_{i \in J} f_i) = [x^k](\prod_{i \in M} f_i)$),
and $[x^0]g = 0$, then for all $J \supseteq M$,
$[x^n](\prod_{i \in J} (f_i \circ g)) = [x^n](\prod_{i \in M} (f_i \circ g))$.
\end{lemma}

\begin{proof}
\leanok
By Lemma~\ref{lem.fps.subs.rule-infprod-fin}, $\prod_{i \in J} (f_i \circ g) = (\prod_{i \in J} f_i) \circ g$.
The approximator condition gives $\prod_{i \in J} f_i \overset{x^n}{\equiv} \prod_{i \in M} f_i$.
By Lemma~\ref{lem.fps.infprod.xnEquiv-comp}, composing with $g$ preserves the equivalence.
\end{proof}

\begin{proposition}
\label{prop.fps.subs.rule-infprod-mulable}
\lean{PowerSeries.comp_prod_multipliable}
\leanhelper
\leanok
If $(f_i)_{i \in I}$ is multipliable and $[x^0]g = 0$, then $(f_i \circ g)_{i \in I}$
is multipliable.
\end{proposition}

\begin{proof}
\leanok
For each $n$, combine the approximators for degrees $0, 1, \ldots, n$ into a single
approximator $M$ (their union), then apply Lemma~\ref{lem.fps.infprod.comp-approx}.
\end{proof}

\subsection{Additional supporting lemmas}

\begin{lemma}
\label{lem.fps.prod.irlv.mul-dvd}
\lean{PowerSeries.coeff_zero_of_dvd}
\leanhelper
\leanok
If $u \mid v$ and $[x^m] u = 0$ for all $m \leq n$, then $[x^m] v = 0$ for all $m \leq n$.
\end{lemma}

\begin{proof}
\leanok
Write $v = u \cdot w$. Then $[x^m] v = \sum_{(i,j) \in \mathrm{antidiag}(m)} [x^i] u \cdot [x^j] w$.
Since $i \leq m \leq n$, we have $[x^i] u = 0$, so each term vanishes.
\end{proof}
